\chapter[Reproducibility and Cost]{Reproducibility, Environment\texorpdfstring{\\}{ }and Cost}
\label{app:reproducibility}

\section{Execution Environments}

The work was not carried out in a single environment, and this appendix records
the three that were used rather than presenting a tidy composite. Each experiment
directory contains a \texttt{manifest.json} recording the interpreter and library
versions actually present when that experiment ran, so any run can be attributed
to its environment without relying on Table~\ref{tab:environments}.

\begin{table}[htbp]
  \centering
  \caption{Environments used, as recorded in the experiment manifests.}
  \label{tab:environments}
  \small
  \begin{tabularx}{\textwidth}{@{}l l l X@{}}
    \toprule
    \textbf{Env.} & \textbf{Python} & \textbf{OpenAI client} & \textbf{Used for} \\
    \midrule
    A & 3.10.11 & 2.24.0 &
      Candidate retrieval runs, both candidate dense-embedding runs, and the
      held-out generation run \\
    B & 3.10.16 (Anaconda) & 2.53.0 &
      Reviewed retrieval runs, both reviewed dense-embedding runs, and the
      development generation runs \\
    C & 3.10.12 (Linux) & --- &
      Context assembly, escalation calibration, and the re-scoring of the
      historical baselines. No hosted call is made by any of these. \\
    \bottomrule
  \end{tabularx}
\end{table}

Table~\ref{tab:environment} details environment~B, the machine on which most of the work was performed.

\begin{table}[htbp]
  \centering
  \caption{Hardware and principal library versions for environment B, the machine on which most of the work was performed.}
  \label{tab:environment}
  \small
  \begin{tabular}{ll}
    \toprule
    \textbf{Component} & \textbf{Value} \\
    \midrule
    Operating system & Windows 10.0.26200 \\
    Processor & Intel64 Family 6 Model 186 Stepping 2 \\
    Logical CPUs & 16 \\
    GPU & NVIDIA GeForce RTX 4070 Laptop \\
    Python & 3.10.16 (CPython, Anaconda) \\
    \midrule
    \texttt{openai} & 2.53.0 \\
    \texttt{numpy} & 2.2.6 \\
    \texttt{torch} & 2.7.1+cu118 \\
    \texttt{transformers} & 5.12.1 \\
    \texttt{matplotlib} & 3.10.0 \\
    \texttt{python-dotenv} & 1.0.1 \\
    \texttt{scikit-learn} & 1.6.0 \\
    \texttt{peft} & 0.12.0 \\
    \texttt{accelerate} & 1.14.0 \\
    \texttt{bitsandbytes} & 0.49.2 \\
    \texttt{safetensors} & 0.5.2 \\
    \bottomrule
  \end{tabular}
\end{table}

\paragraph{A difference between the two generation runs.}
Development generation was performed in environment B and held-out generation in
environment A, so the two runs used different client library versions
(2.53.0 and 2.24.0) and different Python patch releases. This was not intended
and is recorded because Chapter~\ref{chap:results} compares those two runs
directly. The request is determined by the prompt, the schema, the model
identifier and the temperature, all of which are pinned and hashed identically in
both manifests, and the model itself
(\texttt{gpt-4o-mini-2024-07-18}) executes server-side; a client-version effect
is therefore unlikely. It is not excluded. The runs were not repeated to test it:
responses are cached by model, prompt, schema and context, so a repeat would
return the stored answers without contacting the API, and bypassing the cache
would spend the held-out partition a second time, which decision D043 forbids.

\section{Determinism and Caching}

Generation used temperature 0 throughout. Every response is cached on disk under
a key covering the model, the system instruction, the user message and the schema
digest, so a repeated run reproduces the stored answers exactly and cannot be
re-billed. The corpus embeddings were copied between cache directories rather
than recomputed when the benchmark was frozen, because the corpus text did not
change; the copies are recorded in the usage tables with
\texttt{billed = no} so the saving is visible rather than silent.

The evaluation split is deterministic and derived from a string seed consumed by
SHA-256 rather than from a random number generator. The frozen split is
\emph{inherited} from the candidate split rather than recomputed, for the reason
given in decision D028.

\section{Hosted API Usage}

Two model families were used: \texttt{gpt-4o-mini-2024-07-18} for all generation,
and \texttt{text-embedding-3-small} and \texttt{text-embedding-3-large} for dense
retrieval. The LoRA-adapted system was executed locally on the GPU listed above
and made no hosted call.

Table~\ref{tab:usage-completions} reports the completion requests and their tokens.

\begin{table}[htbp]
  \centering
  \caption{Completion requests and tokens, from the provider's usage export. Consolidated in \texttt{docs/billing/}.}
  \label{tab:usage-completions}
  \small
  \begin{tabular}{llrrr}
    \toprule
    \textbf{Date} & \textbf{Run} & \textbf{Requests} & \textbf{Input} & \textbf{Output} \\
    \midrule
    2026-03-18 & Few-shot baseline & 300 & 168{,}466 & 18{,}913 \\
    2026-05-26 & RAG baseline & 300 & 267{,}952 & 17{,}268 \\
    2026-08-09 & E5 generation, development & 84 & 135{,}852 & 28{,}065 \\
    2026-08-12 & E5 generation, held-out & 56 & 90{,}638 & 16{,}462 \\
    \midrule
    & \textbf{Total} & \textbf{740} & \textbf{662{,}908} & \textbf{80{,}708} \\
    \bottomrule
  \end{tabular}
\end{table}

Table~\ref{tab:usage-embeddings} reports embedding tokens separately, deduplicated on the text digest so that copied caches are not counted twice.

\begin{table}[htbp]
  \centering
  \caption{Embedding tokens, from the stored cache metadata, deduplicated on the text digest.}
  \label{tab:usage-embeddings}
  \small
  \begin{tabular}{llrr}
    \toprule
    \textbf{Model} & \textbf{Texts} & \textbf{Tokens} & \textbf{Billed} \\
    \midrule
    \texttt{text-embedding-3-large} & 1{,}535 corpus children & 171{,}186 & yes \\
    \texttt{text-embedding-3-small} & 1{,}535 corpus children & 171{,}186 & yes \\
    \texttt{text-embedding-3-large} & 99 candidate questions & 1{,}635 & yes \\
    \texttt{text-embedding-3-small} & 99 candidate questions & 1{,}635 & yes \\
    \texttt{text-embedding-3-large} & 98 frozen questions & 1{,}736 & yes \\
    \texttt{text-embedding-3-small} & 98 frozen questions & 1{,}736 & yes \\
    \texttt{text-embedding-3-small} & 127 legacy chunks $+$ 98 questions & 58{,}762 & yes \\
    \midrule
    \multicolumn{2}{l}{\textbf{Total billed}} & \textbf{407{,}876} & \\
    \multicolumn{2}{l}{Reused from cache, not re-billed} & 342{,}372 & no \\
    \bottomrule
  \end{tabular}
\end{table}

\section{Cost}

Table~\ref{tab:cost} lists the recorded charges.

\begin{table}[htbp]
  \centering
  \caption{Recorded charges, from the provider's cost export.}
  \label{tab:cost}
  \small
  \begin{tabular}{llr}
    \toprule
    \textbf{Date} & \textbf{What} & \textbf{USD} \\
    \midrule
    2026-03-18 & Few-shot baseline generation & 0.0366 \\
    2026-05-26 & RAG baseline generation & 0.0506 \\
    2026-07-20 & Corpus and question embeddings, both models & 0.0259 \\
    2026-08-08 & Question re-embedding after the freeze & 0.0003 \\
    2026-08-09 & E5 generation, development & 0.0372 \\
    2026-08-12 & E5 generation, held-out & 0.0235 \\
    2026-08-30 & E5 generation, enriched context, development & 0.0276 \\
    2026-08-30 & Few-shot and RAG baselines regenerated over 98 records & 0.0860 \\
    \midrule
    & \textbf{Total} & \textbf{0.2877} \\
    \bottomrule
  \end{tabular}
\end{table}

\textbf{The entire project consumed roughly twenty-nine United States cents of
hosted inference.} The first six charges are taken from the provider's billing
record. The two of 30 August are computed from the token counts the runs
themselves report, at the same unit prices; they were incurred after the last
billing export used here and are marked as such rather than presented as
reconciled.

As a consistency check, five of the first six charges reconstruct exactly from the
recorded token counts at unit prices of \$0.15 and \$0.60 per million tokens for
input and output on \texttt{gpt-4o-mini}, and \$0.02 and \$0.13 per million input
tokens for the small and large embedding models. The remaining charge, for the
RAG baseline on 26 May, exceeds its reconstruction by \$0.00008686; that script
also embeds each query before retrieving, and those embedding calls do not appear
in the completions export.

The cost is worth stating plainly because it bears on the argument of this work.
The expensive parts of building a defensible system here were the corpus repair,
the benchmark construction, and the independent review --- none of which is
inference. Nothing in Chapter~\ref{chap:results} would have been improved by
spending more on model calls, and the two most useful findings, the
citation-legibility defect and the reviewers' disagreement, cost nothing at all
because they came from asking practitioners rather than from running a model.

\section{Reproducing the Results}

\noindent\textbf{Repository.} \url{https://github.com/MoSaleh47/Where-RAG-Breaks---Citation-Grounded-Question-Answering-and-Escalation-over-EU-Aviation-Regulations}

\noindent The repository carries the corpus builders, the retrieval, context and generation drivers, the citation evaluator, the verification layer, the test suite, the decision log and the generated-artefact log. It does not carry the reviewer identity files: the mapping from the pseudonyms R1--R4 to named people is held outside version control and is excluded by \texttt{.gitignore}, so it has never entered the repository's history. Provider billing exports and credentials are excluded on the same basis.

Every experiment writes an immutable directory containing its outputs, a
\texttt{manifest.json} with SHA-256 digests of all inputs and outputs, and the
software versions in use. The order of reconstruction is:

\begin{enumerate}
  \item \texttt{build\_parent\_child\_corpus.py} --- deterministic corpus; the
    rebuild is byte-identical.
  \item \texttt{map\_legacy\_sources.py}, \texttt{build\_legacy100\_ai\_audit.py}
    --- lineage from the historical flat chunks and the evidence-linked
    candidate benchmark.
  \item \texttt{freeze\_reviewed\_benchmark.py} with \texttt{--inherit-splits}
    --- the frozen benchmark. Refuses incomplete or malformed review input.
  \item \texttt{rerun\_reviewed\_retrieval.py} --- all five retrieval
    configurations; resumes from completed stages and reuses cached vectors.
  \item \texttt{rescore\_legacy\_baselines.py} --- historical baselines under the
    corrected evaluator. Makes no model call.
  \item \texttt{build\_candidate\_contexts.py}, \texttt{run\_generation.py} ---
    evidence contexts and generation. The latter defaults to the development
    partition and refuses the held-out partition without an explicit flag.
  \item \texttt{calibrate\_escalation.py} --- escalation signals and rules.
    Fits nothing; makes no model call.
  \item \texttt{capture\_environment.py} --- regenerates the tables in this
    appendix from the manifests and cache metadata.
\end{enumerate}

Unit tests covering the corpus parser, the lineage mapping, grouped splitting,
BM25, fusion, reranking, context assembly, the review contract, citation
resolution and answer verification are run with \texttt{python -m pytest tests/}.
